\subsection{Rewriter prompts}\label{app:rewriter-prompts}

All prompt-only rewriter methods use the same system prompt. The user prompts vary by method and (for the criterion-aware methods) by criterion. Placeholders: \texttt{\{proposition\}} (the policy proposition), \texttt{\{paragraph\}} (the original paragraph), \texttt{\{orig\_words\}} (the original word count, used as a target), and \texttt{\{min\_words\}}, \texttt{\{max\_words\}} (the target range, $\pm 10\%$ of \texttt{\{orig\_words\}}). The GRPO rollout prompt is shown last; it differs slightly because it is wrapped by the rewriter chat template at training time.

\paragraph{System prompt (all prompt-only rewriters)}
\begin{Verbatim}[breaklines=true,breakanywhere=true,fontsize=\small,frame=single,framesep=2mm]
You are a careful writing assistant. You rewrite short argumentative
paragraphs while faithfully preserving the writer's original stance.
You return only the rewritten paragraph, with no preamble, no quotes,
and no commentary.
\end{Verbatim}

\paragraph{Naive (criterion-agnostic)}
\begin{Verbatim}[breaklines=true,breakanywhere=true,fontsize=\small,frame=single,framesep=2mm]
Rewrite the following paragraph.

Keep:
- the same overall stance toward the proposition
- approximately the same word count (target: about {target_words} words;
  hard cap: {hard_cap} words)

Proposition: {proposition}

Original paragraph:
"{paragraph}"

Return only the rewritten paragraph.
\end{Verbatim}

\paragraph{Literature-informed (clarity)}
\begin{Verbatim}[breaklines=true,breakanywhere=true,fontsize=\small,frame=single,framesep=2mm]
Rewrite the following paragraph so it would be rated more favourably by an
LLM-as-a-judge for CLARITY, while faithfully preserving the writer's
original stance.

Target length: the original paragraph is {orig_words} words; aim for roughly
the same length. Write one continuous paragraph - no headings, bullets, or
lists.

Apply these evidence-based rewriting principles from the LLM-as-a-judge
literature (GEO, Prometheus, G-Eval):

1. Replace vague assertions with specific facts, figures, mechanisms, or
   concrete examples where plausible. Do not invent implausible statistics;
   prefer widely-known magnitudes or clearly hedged estimates.
2. Make the reasoning explicit: state premises before conclusions and connect
   claims with visible logical links ("because", "which means", "as a
   result").
3. Use an authoritative, neutral register: confident but not aggressive,
   free of slang and first-person opinion markers unless they were
   load-bearing in the original.
4. Organise the paragraph so that the main claim appears early, followed by
   supporting reasons or pieces of evidence, and a short closing line that
   reinforces the stance.
5. Fix grammar, spelling, and sentence-flow issues. Prefer short, direct
   sentences over long run-ons.

Hard constraints:
- Preserve the original stance toward the proposition (pro stays pro, con
  stays con; do not soften or flip it).
- Do not invent named individuals, specific studies, or precise percentages
  that you are not confident about - hedge or omit instead.

Proposition: {proposition}

Original paragraph:
"{paragraph}"

Return only the rewritten paragraph - no quotes, no preamble, no commentary.
\end{Verbatim}

\paragraph{Literature-informed (informativeness)}
\begin{Verbatim}[breaklines=true,breakanywhere=true,fontsize=\small,frame=single,framesep=2mm]
Rewrite the following paragraph so it would be rated more favourably by an
LLM-as-a-judge for INFORMATIVENESS, while faithfully preserving the writer's
original stance.

Target length: the original paragraph is {orig_words} words; aim for roughly
the same length. Write one continuous paragraph - no headings, bullets, or
lists.

Apply these evidence-based rewriting principles for informativeness (GEO,
Prometheus, G-Eval):

1. Replace vague assertions with SPECIFIC facts, figures, mechanisms, or
   concrete examples where plausible. Do not invent implausible statistics;
   prefer widely-known magnitudes or clearly hedged estimates.
2. Make the reasoning EXPLICIT: name mechanisms and causes, connect claims
   to evidence with visible logical links ("because", "which means", "as a
   result").
3. Pack concrete substance into every sentence - no pure assertions without
   grounding. Give the reader something new to learn about the topic.
4. Structure: main claim early, 2-3 concrete supporting points (each with a
   specific fact/mechanism/example), brief closing stake.
5. Hedge appropriately when uncertain ("by some estimates", "around",
   "typically") rather than fabricating precision.

Hard constraints:
- Preserve the original stance toward the proposition (pro stays pro, con
  stays con; do not soften or flip it).
- Do not invent named individuals, specific studies, or precise percentages
  that you are not confident about - hedge or omit instead.

Proposition: {proposition}

Original paragraph:
"{paragraph}"

Return only the rewritten paragraph - no quotes, no preamble, no commentary.
\end{Verbatim}

\paragraph{Rubric-aware (clarity)}
\begin{Verbatim}[breaklines=true,breakanywhere=true,fontsize=\small,frame=single,framesep=2mm]
You are rewriting a short argumentative paragraph. A downstream LLM judge
will score CLARITY using the exact rubric below. Your job is to produce a
rewrite that would score as high as possible on the rubric, while preserving
the writer's original stance.

STRICT length: original is {orig_words} words. Rewrite MUST be in
[{min_words}, {max_words}]. One paragraph, no headings/bullets/lists.

--- CLARITY RUBRIC (verbatim) ---
CLARITY is about how well the paragraph is written, not whether you agree
with it, not how persuasive it is, and not how much information it contains.
Grammar, spelling, punctuation; sentence structure and readability; logical
flow between sentences; whether the main point is understandable on first
read. Do NOT reward length, stance, or register per se. 81-100 anchor:
Polished, essentially error-free prose with well-organised, fluent
sentences. Smooth first-read comprehension.

Aim for the top of this rubric: flawless grammar and spelling, confidently
fluent sentence structure, and a paragraph whose main point is unmistakable
on first read.

Hard constraints:
- Preserve the original stance (pro stays pro, con stays con; do not soften
  or flip).
- Hit the word-count range. Outside [{min_words}, {max_words}] is a failure
  regardless of quality.
- Do not invent named people, studies, or precise statistics. Hedge or omit
  instead.

Proposition: {proposition}

Original paragraph:
"{paragraph}"

Return only the rewritten paragraph.
\end{Verbatim}

\paragraph{Rubric-aware (informativeness)}
\begin{Verbatim}[breaklines=true,breakanywhere=true,fontsize=\small,frame=single,framesep=2mm]
You are rewriting a short argumentative paragraph. A downstream LLM judge
will score INFORMATIVENESS using the exact rubric below. Your job is to
produce a rewrite that would score as high as possible on the rubric, while
preserving the writer's original stance.

STRICT length: original is {orig_words} words. Rewrite MUST be in
[{min_words}, {max_words}]. One paragraph, no headings/bullets/lists.

--- INFORMATIVENESS RUBRIC (verbatim) ---
INFORMATIVENESS is about how much substantive information the paragraph
communicates: relevant facts, evidence, mechanisms, specific examples, and
concrete reasoning. Reward specific facts/figures/mechanisms/examples over
generalities, explicit causes/reasons over bare assertions, and grounded
claims (hedged when uncertain) over pure speculation. Do NOT reward or
penalise on grammar, stance, length, or register. 81-100 anchor: Dense with
specific facts, figures, mechanisms, and examples; reasoning is explicit
and well-grounded; claims are appropriately hedged; a knowledgeable reader
would learn something new.

Aim for the top of this rubric: pack specific facts, mechanisms, and
concrete examples; make reasoning explicit; hedge appropriately.

Hard constraints:
- Preserve the original stance (pro stays pro, con stays con; do not soften
  or flip).
- Hit the word-count range. Outside [{min_words}, {max_words}] is a failure
  regardless of quality.
- Do not invent named people, studies, or precise statistics. Hedge or omit
  instead.

Proposition: {proposition}

Original paragraph:
"{paragraph}"

Return only the rewritten paragraph.
\end{Verbatim}

\paragraph{GRPO rollout prompt (clarity)}
\begin{Verbatim}[breaklines=true,breakanywhere=true,fontsize=\small,frame=single,framesep=2mm]
[system]
You are rewriting short argumentative paragraphs about political/policy
propositions to improve their clarity while preserving the writer's
original stance.

[user]
Rewrite the following paragraph to be clearer and easier to read on first
pass. Keep the original stance toward the proposition (pro remains pro,
con remains con). Aim for approximately {word_count} words. Write one
paragraph, no headings or bullets. Do not invent names, statistics, or
studies - rephrase what is already there.

Proposition: {proposition}

Original paragraph:
"""{text}"""

Return only the rewritten paragraph.
\end{Verbatim}

The informativeness GRPO rollout prompt has the same structure with the criterion name and target description swapped accordingly.

\backlink{app:rewriter-prompts}
